\documentclass{article}

% Incluir el archivo de variables
%Archivo de variables
\newcommand{\tfg}{Trabajo de Fin de Grado}
\newcommand{\myTfgDate}{6 de diciembre de 2024}
\newcommand{\myAuthorFullName}{Raúl Pradanas Martín}
\newcommand{\myTeacherFullName}{Adrián Domínguez Díaz}
\newcommand{\myUniversityTitle}{Elaboración de una inteligencia artificial capaz de jugar Starcraft 2}
\newcommand{\myUniversityTitleEnglish}{Development of an artificial intelligence capable of playing Starcraft 2}
\newcommand{\myDepartment}{Departamento de informática}


\usepackage[utf8]{inputenc}


\title{Anteproyecto de \tfg}
\date{\myTfgDate}
\author{\myAuthorFullName}


\begin{document}

\maketitle

\begin{description}
    \item [Titulación:] \myUniversityTitle
    \item [Titulacion en inglés:] \myUniversityTitleEnglish
    \item [Tutor:] \myTeacherFullName
    \item [Departamento:] \myDepartment
\end{description}


\section{Introducción}
\label{sec:introduccion}

\textit{En este apartado se describirá el contexto en el que se desenvolverá el
  TFG y sus antecedentes si existen. Recuerda que en el anteproyecto
  también se pueden poner referencias bibliográficas como \cite{moore2002}}.

  \section{Objetivos}
\label{sec:objetivos-y-campo}

\textit{En este apartado se delimitan y explican con claridad los objetivos
  generales a conseguir con el TFG y la aplicación del mismo, así como los
  objetivos específicos en su caso. Algo del tipo:}

El objetivo fundamental de este proyecto es el diseño, implementación y
evaluación de \ldots

Los objetivos específicos de este proyecto son los siguientes:

\begin{itemize}
\item Realizar un  \ldots
\item Diseñar, implementar y evaluar  \ldots, siguiendo la arquitectura
  mostrada en la figura \ref{fig_arquitectura}, y que tendrá las
  siguientes características:

  \begin{enumerate}
  \item Característica 1 \ldots
  \item Característica 2 \ldots
  \item  \ldots
  \item Característica n \ldots
  \end{enumerate}

\item Analizar  \ldots

\item Documentar  \ldots

\end{itemize}

\begin{figure}[tphb]
  \centering
  \caption{Arquitectura del sistema completo.}
  \label{fig_arquitectura}
\end{figure}


\section{Metodología y plan de trabajo}
\label{sec:metodologia-y-plan}

\textit{Aquí se incluirá una descripción (puede ser incluso una enumeración)
  clara de las etapas que se van a seguir, y si es posible se deberá
  incluir un diagrama de Gantt. Por ejemplo algo del estilo a:}

Estas son las fases de desarrollo que se van a seguir para la
consecución de los objetivos del proyecto descritos en la sección~\ref{sec:objetivos-y-campo}:

\begin{enumerate}
  
\item Formación inicial (1 mes)
  
  \begin{itemize}
  \item Formación en \ldots
  \item Consulta de la API \ldots
  \item Consulta bibliográfica \ldots
  \item Profundización en herramientas de soporte \ldots
  \item  \ldots
  \end{itemize}

\item Diseño del entorno \ldots (0,5 meses)

\item Diseño, implementación y evaluación del  \ldots (2 meses)
  \begin{itemize}
  \item Definición del  \ldots
  \item Implementación del  \ldots
  \item Evaluación del  \ldots
  \end{itemize}
  
\item Implementación y evaluación de  \ldots (1 mes)

\item Diseño, implementación y evaluación del módulo  \ldots (1 mes):
  \begin{itemize}
  \item Definición de  \ldots
  \item Definición de  \ldots
  \item Implementación y evaluación  \ldots
  \end{itemize}

\item Integración y  \ldots (1 mes)

\item Documentación  \ldots (0,5 meses)

\end{enumerate}


\section{Medios}
\label{sec:medios}

\textit{Aquí se describen de los medios necesarios para realizar el TFG. Por
  ejemplo:}

Las herramientas que van a ser necesarias para desarrollar este proyecto
son las siguientes:

\begin{itemize}
\item PC compatible
\item Sensor  \ldots
\item Sistema operativo GNU/Linux~\cite{gnulinux}
\item Entorno de desarrollo Emacs/Vim~\cite{emacs}
\item Procesador de textos \LaTeX~\cite{lamport94}
\item Control de versiones CVS~\cite{cvs}
\item Compilador C/C++ gcc~\cite{gcc}
\item Gestor de compilaciones make~\cite{make}
\item Robot con movilidad.
\item  \ldots
\end{itemize}



Otros recursos necesarios para la elaboración del proyecto son:

\begin{itemize}
\item Herramientas  \ldots
\item Sistema de desarrollo  \ldots
\item  \ldots
\end{itemize}

\end{document}